\documentclass[12pt, letterpaper]{article}

% ====================================================================
% PAQUETES BÁSICOS Y CONFIGURACIÓN DEL IDIOMA
% ====================================================================

\usepackage[spanish,es-tabla]{babel} % Configura el documento en español y usa "Tabla" en vez de "Cuadro"
\usepackage{float} % Para el especificador [H]
\usepackage[utf8]{inputenc}   % Permite escribir acentos y eñes directamente
\usepackage[spanish,es-tabla]{babel} % Configura el documento en español y usa "Tabla" en vez de "Cuadro"
\usepackage[margin=2.5cm]{geometry} % Fija los márgenes de las páginas a 2,5 cm
\usepackage{graphicx}         % Permite insertar imágenes y diagramas
\usepackage{booktabs}         % Para hacer tablas limpias y profesionales
\usepackage{tabularx}
\usepackage{amsmath, amssymb} % Paquetes matemáticos estándar
\usepackage{microtype}        % Mejoras en la tipografía general y espaciado
\usepackage{hyperref}         % Crea enlaces clicables en el índice
\hypersetup{
    colorlinks=true,
    linkcolor=blue,
    filecolor=magenta,      
    urlcolor=cyan,
    pdftitle={Proyecto Banda Transportadora Desviadora de Metales},
}

% ====================================================================
% DATOS DE LA PORTADA
% ====================================================================
\title{
    \vspace{2cm}
    \Huge \textsf{\textbf{Banda Transportadora Desviadora de Metales}} \\
    \vspace{1cm}
    \large Materia: Estructuras Computacionales \\
    \vspace{2cm}
}

\author{
    \textbf{Juan Felipe Pachón Restrepo} \\
    C.C. 1054542918 \\
    \vspace{1.5cm} \\
    Profesor: \\
    \textbf{Gustavo Adolfo Osorio Londoño}
}

\date{
    \vfill
    Universidad Nacional de Colombia \\
    Sede Manizales \\
    Campus La Nubia
}

% ====================================================================
% CUERPO DEL DOCUMENTO
% ====================================================================
\begin{document}

% 1. Generación automática de la portada
\maketitle
\thispagestyle{empty} % Quita el número de página en la portada
\clearpage

% 2. Índice automático del reporte
\tableofcontents
\clearpage

\section{Proyecto / Problema / Sistema}
\label{sec:proyecto_problema_sistema}

El sistema recibe información desde un sensor inductivo encargado de detectar piezas metálicas y desde un botón de usuario configurado para identificar múltiples pulsaciones o presiones prolongadas. Estos datos son procesados por una máquina de estados que se ejecuta de forma asíncrona sobre un sistema operativo en tiempo real, generando una respuesta observable en el motor de tracción de la banda, la posición angular de un servomotor de desvío y un indicador LED de estado.

El prototipo interactúa directamente con su entorno físico al identificar materiales ferrosos sobre la cinta transportadora en movimiento y sin necesidad de contacto directo. Al ocurrir una detección, el brazo deflector altera la trayectoria del objeto para retirarlo de la línea antes de regresar a su posición inicial de reposo. Asimismo, el operador puede intervenir en la marcha del sistema mediante el botón de control para encender la banda, cambiar su velocidad de manera inmediata, apagarla o activar una detención de emergencia total si es necesario. Todo el comportamiento del problema se define bajo un enfoque puramente funcional y de control con el entorno, prescindiendo del uso de nombres de pines, líneas de código o librerías de software específicas en esta sección introductoria.

\clearpage

\section{Entradas, Salidas y Eventos}
\label{sec:entradas_salidas_eventos}

En esta sección se organizan los componentes de hardware que actúan como periféricos de entrada y salida, así como los eventos lógicos que gobiernan el firmware.

\begin{table}[H]
\centering
\small
\renewcommand{\arraystretch}{1.1}
\caption{Clasificación de entradas, salidas y eventos del sistema.}
\label{tab:entradas_salidas_eventos}
\vspace{0.3cm}
\begin{tabularx}{\textwidth}{lXX}
\toprule
\textbf{Elemento} & \textbf{Descripción en el Prototipo} & \textbf{Tipo de Señal / Conexión} \\
\midrule
Entrada 1         & Sensor inductivo de metales LJ12A3 & Digital (Configurado con interrupción) \\
Entrada 2         & Botón de usuario integrado (PC13) & Digital (Detección de ambos flancos de conmutación) \\
\midrule
Salida 1          & Control de velocidad del motor & Modulación por ancho de pulso (PWM) \\
Salida 2          & Control de posición del brazo (Servo) & Modulación por ancho de pulso (PWM) \\
Salida 3a         & Dirección del motor - Línea IN1 del puente H & Digital de salida \\
Salida 3b         & Dirección del motor - Línea IN2 del puente H & Digital de salida \\
Salida 4          & LED indicador de estado (PA5) & Digital de salida \\
\midrule
Evento 1          & Pulsación única o doble en el botón & Entrada por interrupción para control de marcha \\
Evento 2          & Presión prolongada del botón (2 segundos) & Disparo inmediato de la parada de emergencia \\
Evento 3          & Pulsación del botón estando en emergencia & Evento de desbloqueo del sistema \\
Evento 4          & Detección de material metálico en la cinta & Entrada asíncrona (Lectura de 0V en el pin) \\
Evento 5          & Fin del tiempo de retención del brazo & Temporizador controlado por el kernel \\
\bottomrule
\end{tabularx}
\end{table}

\section{Requerimientos y Validation}
\label{sec:requerimientos_preliminares}

A continuación se detallan las especificaciones obligatorias del diseño. Cada requerimiento se define de forma observable y verificable, vinculándose directamente con un criterio explícito para su validación.

\subsection{REQ-01: Gestión de Concurrencia}
El sistema debe administrar la lógica de control, la lectura de sensores y la actualización de actuadores utilizando las herramientas de hilos y tareas del sistema operativo en tiempo real.

\vspace{0.5cm}

\noindent \textbf{Criterio de aceptación:} \\
Se considera cumplido si la máquina de estados procesa los eventos mediante colas de trabajo (\textit{workqueues}) e hilos independientes, evitando retrasos bloqueantes en el flujo principal.

\subsection{REQ-02: Respuesta Asíncrona}
El firmware debe reaccionar de manera inmediata ante los estímulos externos provocados por el botón de control o el sensor de metales.

\vspace{0.5cm}

\noindent \textbf{Criterio de aceptación:} \\
Se considera cumplido si el cambio eléctrico en los pines de entrada dispara una rutina de servicio de interrupción (ISR) por hardware, eliminando el uso de lecturas cíclicas por software (\textit{polling}).

\subsection{REQ-03: Control de Velocidad por Pasos}
El motor de tracción debe iniciar completamente detenido y aumentar su velocidad de forma progresiva en 4 pasos fijos controlados mediante el botón de usuario.

\vspace{0.5cm}

\noindent \textbf{Criterio de aceptación:} \\
Se considera cumplido si el valor inicial de software del $60\%$ se traduce en un estado de reposo ($0\%$ de movimiento real en la banda) y cada pulsación única incrementa el ciclo de trabajo en pasos del $10\%$ hasta alcanzar el máximo ($100\%$) en el cuarto paso.

\subsection{REQ-04: Activación del Brazo Clasificador}
El servomotor encargado del desvío de piezas debe activarse únicamente cuando el sensor inductivo confirme la presencia física de material metálico en la banda.

\vspace{0.5cm}

\noindent \textbf{Criterio de aceptación:} \\
Se considera cumplido si el brazo deflector cambia de su ángulo de reposo al de expulsión si y solo si la máquina de estados procesa el evento de detección de metal y el sistema no se encuentra apagado.

\subsection{REQ-05: Temporización del Rechazo}
El brazo mecánico debe permanecer extendido el tiempo necesario para asegurar que el objeto sea retirado por completo de la trayectoria de la banda antes de regresar a la posición inicial.

\vspace{0.5cm}

\noindent \textbf{Criterio de aceptación:} \\
Se considera cumplido si el temporizador del kernel retiene el servomotor en la posición de desvío durante exactamente 4,0 segundos antes de ordenar automáticamente su retorno.

\subsection{REQ-06: Inicialización Segura}
Al energizarse o reiniciarse la tarjeta de desarrollo, el sistema debe iniciar obligatoriamente en una condición segura y pasiva para evitar movimientos accidentales.

\vspace{0.5cm}

\noindent \textbf{Criterio de aceptación:} \\
Se considera cumplido si al arrancar el microcontrolador el motor se configura en un $0\%$ de marcha real, el servo se posiciona de forma síncrona en resguardo ($0^\circ$) y la FSM se sitúa firmemente en el estado de apagado.

\subsection{REQ-07: Protección del Puente H}
El código de control del motor debe impedir cualquier combinación lógica de salida que ponga en riesgo la integridad física de la etapa de potencia.

\vspace{0.5cm}

\noindent \textbf{Criterio de aceptación:} \\
Se considera cumplido si las funciones de la API aseguran que las líneas digitales de dirección IN1 e IN2 nunca compartan un valor lógico alto (1) de forma simultánea bajo ninguna circunstancia.

\subsection{REQ-08: Parada de Emergencia Absoluta}
El sistema debe contar con un mecanismo de interrupción drástica que detenga de inmediato toda actividad física ante condiciones imprevistas.

\vspace{0.5cm}

\noindent \textbf{Criterio de aceptación:} \\
Se considera cumplido si mantener presionado el botón por 2,0 segundos apaga el motor, retorna el servo a posición base de manera inmediata y bloquea cualquier nueva detección del sensor mediante un filtro absoluto en la entrada.

\section{Estados del Sistema}
\label{sec:estados_preliminares}

El comportamiento secuencial del programa se organiza a través de cuatro estados lógicos principales:

\begin{itemize}
    \item \textbf{\texttt{STATE\_SYSTEM\_OFF}:} Estado inicial de seguridad. La banda transportadora permanece detenida (el ciclo de trabajo PWM está configurado en su límite mínimo del $60\%$, lo que equivale al $0\%$ de movimiento real) y el servomotor se encuentra firmemente asegurado en su ángulo de reposo mediante una aduana de software que anula lecturas tardías del sensor.
    \item \textbf{\texttt{STATE\_BANDA\_RUNNING}:} Estado de operación normal. El motor de tracción desplaza la cinta de manera continua a la velocidad seleccionada por el usuario (ajustable en pasos del $10\%$), mientras el brazo se mantiene estático a la espera de metal.
    \item \textbf{\texttt{STATE\_METAL\_REJECT}:} Estado de seguridad y desvío. Se activa al instante cuando el sensor detecta una pieza metálica. El servomotor se extiende a $60^\circ$ para retirar el objeto de la pista durante un tiempo fijo de 4,0 segundos antes de restaurar la marcha normal.
    \item \textbf{\texttt{STATE\_EMERGENCY\_STOP}:} Estado de bloqueo crítico. Detiene por completo la potencia del motor, devuelve el servo a la posición base y activa un parpadeo visual rápido en el LED de estado ($100\,$ms). En este estado, el sensor inductivo queda completamente sordo por software. Un toque simple en el botón desbloquea este estado, enviando el sistema directamente a la marcha operativa mínima.
\end{itemize}

\section{Eventos del Programa}
\label{sec:eventos_preliminares}

Los eventos definidos en el firmware controlan las transiciones de la máquina de estados:

\begin{itemize}
    \item \textbf{\texttt{EVENT\_CLICK\_INCREMENT}:} Se genera al detectar un único clic válido en el botón dentro de la ventana de tiempo. Si el sistema está apagado, cambia al estado de marcha. Si la banda ya está corriendo, incrementa la velocidad del motor.
    
    \item \textbf{\texttt{EVENT\_CLICK\_DECREMENT}:} Se genera al registrarse un doble clic válido dentro de la ventana de tiempo. Provoca una reducción progresiva de la velocidad del motor en pasos del $10\%$. Si el sistema ya se encuentra en su velocidad mínima operativa ($70\%$), este evento disminuye el valor al $60\%$, lo que apaga físicamente el motor de tracción, cancela los temporizadores activos del brazo y regresa de forma segura al estado \texttt{STATE\_SYSTEM\_OFF}.
    
    \item \textbf{\texttt{EVENT\_BUTTON\_PRESSED}:} Se dispara al instante en que el usuario presiona el botón físico. Actúa como señal de entrada primaria para evaluar tiempos prolongados o como la señal de desbloqueo inmediato cuando el sistema está en emergencia, forzando la transición directa hacia el estado de marcha con velocidad mínima.
    
    \item \textbf{\texttt{EVENT\_EMERGENCY}:} Se gatilla tras mantener presionado el botón de usuario de manera continua por 2000\,ms, forzando al sistema a entrar en la rutina de bloqueo de seguridad.
    
    \item \textbf{\texttt{EVENT\_METAL\_DETECTED}:} Se dispara de forma asíncrona mediante la interrupción del sensor inductivo al detectar metal (conmutación de la señal a 0V). Fuerza la transición inmediata hacia el estado \texttt{STATE\_METAL\_REJECT}.
    
    \item \textbf{\texttt{EVENT\_TIMEOUT}:} Generado automáticamente por el sistema operativo al concluir el periodo de retención de 4000\,ms. Restaura el servo a su posición inicial y regresa la FSM al estado de marcha regular.
\end{itemize}

\clearpage

\section{Diagrama Máquina de Estados del Sistema}
\label{sec:diagrama_fsm}

En esta sección se presenta la Máquina de Estados Finita (FSM) que reacciona de manera inmediata a los eventos del usuario.

\begin{figure}[H]
    \centering    \includegraphics[width=1\textwidth]{maquina.png}
    \caption{Diagrama de la máquina de estados finita (FSM)}
    \label{fig:maquina_estados}
\end{figure}

\section{Hardware y Materiales Disponibles}
\label{sec:hardware_materiales}

A continuación se presenta el inventario consolidado de los componentes utilizados en la implementación del prototipo:

\vspace{0.3cm}

\begin{table}[H]
\centering
\renewcommand{\arraystretch}{1.2}
\caption{Componentes de hardware utilizados en el proyecto.}
\label{tab:hardware_materiales}
\vspace{0.3cm}
\begin{tabularx}{\textwidth}{lX}
\toprule
\textbf{Elemento} & \textbf{Función y Observaciones dentro del Diseño} \\
\midrule
STM32L4 / NUCLEO-L476RG & Tarjeta de desarrollo principal encargada de ejecutar el firmware en Zephyr RTOS. \\
\midrule
Sensor Inductivo LJ12A3 & Sensor industrial de tipo NPN colector abierto, configurado con resistencia de pull-up. \\
\midrule
Puente H L298N & Módulo de potencia empleado para controlar la velocidad (PWM) y sentido del motor DC. \\
\midrule
Servomotor SG90 & Actuador de posición controlado por PWM (50\,Hz) para mover el brazo deflector. \\
\midrule
Chasis y Banda Transportadora & Mecanismo físico de desplazamiento acoplado al motor de tracción continua. \\
\midrule
Fuente de Alimentación Externa & Suministro eléctrico independiente de 12v. \\
\bottomrule
\end{tabularx}
\end{table}

\clearpage

\section{Pines y Periféricos de Conexión}
\label{sec:pines_perifericos_tentativos}

Asignación física y distribución de señales sobre los pines del microcontrolador de la placa Nucleo-L476RG:

\vspace{0.3cm}

\begin{table}[H]
\centering
\renewcommand{\arraystretch}{1.2}
\caption{Distribución de pines y periféricos de hardware.}
\label{tab:pines_perifericos}
\vspace{0.3cm}
\begin{tabularx}{\textwidth}{lXXX}
\toprule
\textbf{Señal de Control} & \textbf{Pin Nucleo} & \textbf{Periférico Interno} & \textbf{Aplicación en el Sistema} \\
\midrule
\texttt{ENA (PWM\_MOTOR)} & PA0 & TIM2 (CH1) & Modulación PWM de la velocidad de la banda (10\,kHz). \\
\midrule
\texttt{IN1 (DIR\_MOTOR\_1)} & PB0 & GPIO & Señal digital de salida para dirección del motor DC. \\
\midrule
\texttt{IN2 (DIR\_MOTOR\_2)} & PB3 & GPIO & Señal digital de salida para dirección del motor DC. \\
\midrule
\texttt{PWM\_SERVO} & PB4 & TIM3 (CH1) & Control de posición angular para el servo SG90 (50\,Hz). \\
\midrule
\texttt{SENSOR\_IN} & PA1 & GPIO & Entrada digital con interrupción por flanco de bajada. \\
\midrule
\texttt{BUTTON\_USER} & PC13 & GPIO & Botón azul de la placa integrado para la captura de clics. \\
\midrule
\texttt{LED\_USER} & PA5 & GPIO & LED verde integrado para indicación visual de los estados. \\
\bottomrule
\end{tabularx}
\end{table}

\clearpage

\section{Arquitectura Modular del Software}
\label{sec:arquitectura_software}

El firmware está estructurado de forma modular en archivos independientes, separando el control de bajo nivel de la lógica de control del sistema:

\begin{table}[H]
\centering
\renewcommand{\arraystretch}{1.2}
\caption{Estructura de archivos de código y responsabilidades.}
\label{tab:arquitectura_software}
\vspace{0.3cm}
\begin{tabularx}{\textwidth}{lXX}
\toprule
\textbf{Módulo / Archivo} & \textbf{Responsabilidad Técnica} & \textbf{Vínculo con el Sistema} \\
\midrule
\texttt{main.c / .h} & Inicialización de periféricos base de la tarjeta, enlace de la ISR del botón y arranque seguro. & Punto de entrada de la aplicación; delega el control a la FSM. \\
\midrule
\texttt{app\_fsm.c / .h} & Implementación de la máquina de estados, control de tareas diferidas, protección de actuadores en reposo y decodificación de clics con filtros contra hilos retrasados. & Núcleo del firmware; procesa eventos y gobierna las APIs de hardware. \\
\midrule
\texttt{lj12a3\_api.c / .h} & Configuración y lectura del sensor inductivo mediante rutinas de interrupción y funciones lógicas de habilitación en caliente. & Notifica de forma asíncrona los eventos de detección metálica bajo permiso de la FSM. \\
\midrule
\texttt{sg90\_api.c / .h} & Controlador PWM para el servomotor. Realiza el cálculo del ancho de pulso angular. & Modifica la posición del brazo extractor ante comandos de la FSM. \\
\midrule
\texttt{l298h\_api.c / .h} & Controlador del Puente H. Gestiona pines de dirección y ciclo de trabajo del motor de tracción. & Modifica la velocidad y parada de la banda transportadora. \\
\bottomrule
\end{tabularx}
\end{table}

\section{Drivers y Servicios del Sistema Operativo}
\label{sec:drivers_servicios}

A continuación se detallan y justifican de manera individual los servicios nativos del kernel de Zephyr RTOS que fueron seleccionados e implementados para dar solución al firmware de la banda transportadora:

\subsection{GPIO (General Purpose Input/Output)}
Este servicio se utiliza para gestionar las líneas digitales básicas del microcontrolador. Específicamente, controla los niveles lógicos altos y bajos de los pines de dirección del puente H (líneas IN1 e IN2) para manejar el sentido de giro del motor DC, y permite el control de encendido o apagado del LED de estado de la tarjeta de desarrollo.

\subsection{GPIO callback / interrupción}
Se emplea para configurar y administrar las interrupciones físicas por hardware en los pines de entrada. Esto permite que el sistema responda de forma asíncrona ante los flancos de conmutación (ambos flancos) del botón de usuario o los flancos de bajada provocados por la activación del sensor inductivo al detectar un objeto metálico, eliminando la necesidad de realizar lecturas cíclicas por software.

\subsection{PWM (Pulse Width Modulation)}
Este controlador se utiliza para generar las señales periódicas de temporización que requieren los actuadores. Se aplica en el motor de tracción para modular el ciclo de trabajo y modificar la velocidad de la cinta transportadora por pasos, y en el servomotor para establecer los anchos de pulso necesarios (a una frecuencia fija de $50\,$Hz) que definen su posición angular exacta.

\subsection{k\_timer / k\_work (Temporizadores y colas de trabajo)}
Se utilizan para planificar tareas diferidas en el tiempo sin bloquear la ejecución del hilo principal de procesamiento. Son indispensables para controlar con precisión los 4,0 segundos que el brazo deflector debe retener su posición durante el rechazo de material, mitigar el ruido eléctrico del botón de usuario mediante un filtrado de antirebote asíncrono por flanco de cola, así como para gestionar de manera independiente la ventana de tiempo para el conteo de clics y el parpadeo del LED en la parada de emergencia.

\subsection{threads (Hilos de ejecución)}
Este servicio permite la ejecución concurrente de tareas bajo el esquema multitarea del sistema operativo en tiempo real. Se utiliza para crear un hilo cooperativo independiente exclusivo para el administrador del motor, separando la lógica de control de la máquina de estados de las actualizaciones de hardware y optimizando el uso de la CPU.

\subsection{k\_sem (Semáforos)}
Se implementan como herramientas de sincronización entre diferentes tareas del kernel. Específicamente, se utiliza un semáforo de control para mantener en estado de espera permanente al hilo encargado del motor, despertándolo de manera inmediata únicamente cuando ocurre un cambio real en la velocidad del sistema tras un evento válido de la FSM.

\section{Evidencia de Validación y Pruebas}
\label{sec:evidencia_validacion}

A continuación se presenta el plan de pruebas sistemático diseñado para verificar el correcto funcionamiento del prototipo. El proceso se estructuró bajo una estrategia de integración incremental, permitiendo validar cada componente de forma aislada antes de evaluar el comportamiento del sistema completo en su conjunto.

\subsection{TEST-01: Prueba independiente del servomotor}
Esta primera etapa se enfoca en verificar la correcta respuesta del actuador de desvío de forma totalmente aislada de los estímulos externos.

\vspace{0.5cm}

\noindent \textbf{Procedimiento del ensayo:} \\

Consiste en invocar de manera directa en el código de arranque la función de control, forzando al servomotor a cambiar su posición angular de forma repetitiva entre los $0^\circ$ (posición de reposo) y los $60^\circ$ (ángulo estimado para la expulsión de piezas).

\vspace{0.5cm}

\noindent \textbf{Resultado esperado:} \\

El brazo deflector debe desplazarse de manera uniforme hacia las posiciones indicadas sin vibraciones excesivas ni saltos abruptos. Asimismo, se debe verificar mediante un osciloscopio que la señal PWM generada por el periférico interno mantenga la frecuencia fija de $50\,$Hz y los anchos de pulso correctos para cada ángulo.

\subsection{TEST-02: Prueba de servomotor con sensor inductivo}
Una vez garantizado el movimiento del actuador, se introduce la primera línea de entrada asíncrona para evaluar el lazo cerrado de detección física.

\vspace{0.5cm}

\noindent \textbf{Procedimiento del ensayo:} \\

Con la banda transportadora completamente detenida en su estado pasivo, se energiza el sensor inductivo LJ12A3 y se acerca manualmente un material ferroso a su zona de detección para observar la reacción inmediata del brazo mecánico.

\vspace{0.5cm}

\noindent \textbf{Resultado esperado:} \\

El sensor debe capturar la presencia del material y conmutar su salida a 0V. Esta transición eléctrica tiene que disparar la rutina de servicio de interrupción (ISR) por hardware, la cual debe despachar con éxito el evento hacia la máquina de estados para que el servomotor se extienda de forma inmediata a la posición de desvío.

\subsection{TEST-03: Prueba de servomotor con sensor y motor en marcha}
En esta fase se integran los dos actuadores principales del diseño junto al sensor industrial para evaluar el comportamiento dinámico del flujo de materiales.

\vspace{0.5cm}

\noindent \textbf{Procedimiento del ensayo:} \\

Se activa la marcha del motor de tracción continua a través del puente H para iniciar el movimiento constante de la cinta transportadora. Posteriormente, se deja rodar una pieza metálica sobre la línea para que pase de forma natural por el frente del área de lectura del sensor inductivo.

\vspace{0.5cm}

\noindent \textbf{Resultado esperado:} \\

La banda debe avanzar de forma suave y continua. Al momento exacto en que la pieza cruza el sensor, el brazo deflector tiene que extenderse para alterar la trayectoria del objeto y retirarlo de la línea. El servomotor debe retener su posición durante los 4,0 segundos programados en el kernel y luego regresar automáticamente a su posición base.

\subsection{TEST-04: Implementación de botón de usuario y LED de estado con el sistema operativo}
En esta etapa se integra formalmente la interfaz de control del operador (botón y LED) junto con el conjunto de hardware validado en los ensayos anteriores (puente H, motor y servomotor).

\vspace{0.5cm}

\noindent \textbf{Procedimiento del ensayo:} \\

Con todo el hardware conectado, se presiona el botón físico PC13 para encender el sistema desde el estado de reposo y se ejecutan pulsaciones para verificar los ciclos de encendido y apagado de la banda. Se observa de forma simultánea que el LED responda visualmente a la interacción del usuario y que el puente H altere la marcha del motor de tracción según las órdenes recibidas.

\vspace{0.5cm}

\noindent \textbf{Resultado esperado:} \\

El sistema debe realizar de forma correcta el encendido y apagado del motor DC a través del puente H de manera coordinada. El LED verde integrado (PA5) debe encenderse al instante con cada clic del usuario como confirmación visual de que la pulsación fue recibida por el microcontrolador. Al apagar la banda mediante la reducción de velocidad al mínimo, el motor debe detenerse por completo y el LED debe quedar apagado, manteniendo el servomotor firmemente resguardado en su posición base de $0^\circ$.

\subsection{TEST-05: Implementación final combinando todo el sistema}
La última etapa evalúa el prototipo en su totalidad, verificando la interacción simultánea de todos los módulos de software, hardware y las directrices de seguridad de la máquina de estados.

\vspace{0.5cm}

\noindent \textbf{Procedimiento del ensayo:} \\

Se pone en marcha el prototipo completo interactuando activamente con el botón para variar la velocidad de la banda transportadora en sus 4 pasos fijos y dejando pasar múltiples piezas (tanto plásticas como metálicas) sobre la cinta en movimiento. Finalmente, se provoca una condición de riesgo simulada y se mantiene presionado el botón por 2,0 segundos para forzar la parada de emergencia.

\vspace{0.5cm}

\noindent \textbf{Resultado esperado:} \\

El motor de la banda debe responder con agilidad a cada clic del operador  o decrementando su velocidad en pasos regulados del $10\%$. El sistema debe ignorar las piezas de plástico y desviar con precisión únicamente los objetos metálicos. Al activar la parada de emergencia, la cinta debe detenerse a cero de forma instantánea. Al pulsar nuevamente el botón para salir de la emergencia, el sistema debe saltar de inmediato al estado de marcha activa, encendiendo el motor directamente en su velocidad mínima operativa ($70\%$).

\end{document}